\section{Introduction}

% ==========================================
% 1. SECTION 1.1
% ==========================================

\subsection{Predictive Modeling in Low-Data Bioactivity Regimes}

Predictive modeling in quantitative structure-activity relationships (QSAR) operates under strict constraints when data is scarce \cite{winkler2026}. Under these conditions, simpler interpretable approaches—such as multivariate linear regression or classical tree-based ensembles—frequently match or exceed the prediction reliability of complex artificial neural networks \cite{hadas2023, xiewt2023}. These parsimonious models maintain higher bias and lower variance. They resist the extreme sensitivity to random fluctuations that plagues overparameterized algorithms, particularly when supplied with robust, quantum-chemically derived descriptors \cite{xiewt2023}.

Small datasets inherently invite overfitting. This bottleneck worsens exponentially when the dimensionality of molecular descriptors exceeds the available sample size \cite{hawkins2004, godarzi2013}. To circumvent data sparsity, contemporary workflows frequently deploy data augmentation (DAU) or synthetic sample generation to artificially expand the training manifold \cite{gerz2026, tyler2025}. However, these interventions demand substantial computational overhead \cite{tyler2025, calle2025}. More critically, improper DAU application obscures the true predictive capacity of the base algorithms. When augmentation precedes dataset partitioning (the ``augment-then-split'' fallacy), it violates independent and identically distributed (i.i.d.) assumptions. Models memorize synthetic artifacts and latent data dependencies rather than learning genuine physicochemical rules, yielding artificially inflated performance metrics \cite{gerz2026, john2025, kannoth2024}.

Evaluating complex architectures in these regimes requires absolute methodological isolation. Hyperparameter tuning on small samples with large parameter spaces inevitably leads to model selection bias \cite{cawley2010}. Consequently, assessing machine learning architectures in isolation provides zero evidence of real-world utility unless the model selection procedure itself undergoes rigorous validation. Double (or nested) cross-validation becomes mandatory to decouple model building from performance estimation \cite{bauman2014, cawley2010}. Standard k-fold cross-validation systematically overestimates explorative capabilities \cite{kyrali2022, xiong2020}, and generative deep learning struggles to extrapolate within sparse, discrete chemical spaces like Simplified Molecular-Input Line-Entry System (SMILES) strings \cite{yifan2023}. To preserve methodological integrity, any generative or augmentative process must remain strictly confined to the training folds after establishing an independent partition \cite{gerz2026, john2025}.

To safely harness overparameterized models \cite{winkler2026} without succumbing to the overfitting trap, empirical safeguards against chance correlation are non-negotiable. Techniques such as y-randomization (response shuffling) explicitly verify that the algorithm maps true chemical relationships rather than memorizing random noise \cite{rucker2007}. Although frequently recommended in literature, its empirical execution is routinely omitted in modern deep learning QSAR studies; enforcing this protocol is therefore essential to definitively rule out chance correlation before advancing candidates to virtual screening. When descriptors undergo meticulous curation and pass these randomization checks, machine learning effectively guides rational design—successfully discriminating potent bioactives, such as Amyloid-$\beta$ aggregation inhibitors, from decoys \cite{avantika2025}. 

Alternatively, transfer learning bypasses data scarcity by adapting base networks optimized on massive source domains \cite{samuel2023}. By enforcing reduced learning rates during fine-tuning \cite{damilola2025} and prioritizing efficient convolutional architectures over heavier Transformer models \cite{riza2025}, researchers can navigate class imbalance and model selection bias in low-resource settings.

% ==========================================
% 1. SECTION 1.2
% ==========================================
\subsection{Domain-Engineered Features versus End-to-End Representation Learning}

Traditional feature-based machine learning relies on domain-engineered descriptors and molecular fingerprints \cite{rajesh2026}. Representations such as Morgan (ECFP4) encode structural topologies, while physicochemical and 3D descriptors capture specific molecular interactions deterministically without requiring massive training corpora \cite{avantika2025, beatrice2026, xiewt2023}. Embedding theoretical chemistry directly into the model ensures high interpretability and algorithmic transparency \cite{xiewt2023}. However, accumulating these handcrafted features generates highly dimensional, redundant matrices \cite{godarzi2013, hadas2023}. In restricted datasets, this redundancy amplifies computational costs and exacerbates the feature selection bottleneck \cite{godarzi2013, hadas2023}.

End-to-end representation learning bypasses manual feature selection. Deep learning paradigms—including Autoencoders (VAEs), Generative Adversarial Networks (GANs), and Graph Neural Networks (GNNs)—autonomously map discrete molecular topologies into continuous latent spaces to model complex, nonlinear interactions \cite{abul2026, rajesh2026, yifan2023}. Similarly, Chemical Language Processing (CLP) applies natural language techniques to string representations like SMILES \cite{riza2025}. Yet, learning representations in discrete chemical spaces remains intrinsically fragile. Altering a single character in a SMILES string can irreparably invalidate the entire molecular structure \cite{yifan2023}.

The utility of deep learning complexity in low-data regimes remains highly contested. These algorithms require dense data manifolds, rendering representation learning ineffective when ground-truth labels are scarce \cite{sotirios2025}. To circumvent this limitation in \textit{Toxoplasma gondii} dihydrofolate reductase (TgDHFR) inhibition, recent studies deploy highly complex pipelines. \cite{karimova2025} employed Gaussian noise-based data augmentation and deep neural network (DNN) ensembles to stabilize training on a restricted dataset, elevating their baseline $R^2$ from \PaperBaselineRTwo{} to \PaperFinalRTwo{}. This intervention prompts a critical evaluation: is synthetic data generation strictly necessary to achieve predictive reliability? Parsimonious classical ensembles frequently match or exceed the predictive fidelity of deep architectures in data-scarce scenarios, avoiding the catastrophic overfitting inherent to unregularized overparameterization \cite{beatrice2026, hadas2023, rajesh2026, salman2023, samuel2023, xiewt2023}.

Traditional skepticism toward nonlinear models in limited-data scenarios is evolving \cite{Dalmau2025}. Properly regularized nonlinear algorithms, hybridized with structured linear priors, achieve competitive results comparable to robust classical methods \cite{Dalmau2025, salman2023}. Transfer learning serves as a primary catalyst for this shift. Fine-tuning pretrained networks like ChemBERTa allows researchers to extract accurate representation learning predictions using only dozens of target data points \cite{latefa2025, samuel2023}.

% ==========================================
% 1. SECTION 1.3
% ==========================================
\subsection{Methodological Rigor and the Threat of Data Leakage in Molecular Modeling}

Data leakage—the unintended introduction of test data information into the training process—is a statistically silent error that fundamentally undermines the validity of machine learning evaluations by inflating metrics without triggering code warnings \cite{brva2026, gerz2026}. In low-data QSAR modeling, this vulnerability frequently manifests across two primary vectors: preprocessing and model selection. When transformations like feature scaling, imputation, or variance-based filtering are applied globally across an entire dataset, the algorithm gains illicit knowledge of the test distribution. This violation of the independent and identically distributed (i.i.d.) assumption yields spuriously high performance estimates \cite{correa2025, lewis2023}.

Similarly, traditional internal validation often suffers from model selection bias. If a validation set continuously influences hyperparameter tuning, it loses its statistical independence. Highly parameterized architectures, such as deep neural networks, are particularly prone to exploiting this overlap. They adapt to noise and produce deceptively optimistic metrics that severely underestimate true generalization errors \cite{bauman2014, sotirios2025}. To preserve methodological integrity, model selection and preprocessing must operate as strictly encapsulated components of the fitting process. Insufficient reporting of these partitioning mechanics often confounds scientific literature, leading researchers to incorrectly attribute high performance to novel complex architectures when the success actually stems from silent data leakage \cite{rojas2026}.

% ==========================================
% 1. SECTION 1.4
% ==========================================

\subsection{Objectives and Principal Contributions}

This study bridges the gap between stringent statistical validation and advanced representation learning to accelerate targeted drug discovery in low-data regimes. Our principal contributions are defined as follows:

\begin{itemize}
	\item \textbf{Methodological Reproducibility and Leak-Free Validation:} Implementation of a fully transparent, reproducible Python pipeline built upon a rigorous $\OuterKFolds \times \OuterRepeats$ repeated nested cross-validation protocol (\TotalOuterFolds{} outer evaluations). This architecture ensures absolute statistical independence between hyperparameter optimization and model evaluation, eradicating model selection bias and data leakage.
	
	\item \textbf{Rigorous Statistical Inference in QSAR:} Introduction of the Nadeau-Bengio variance correction to molecular machine learning evaluations. By explicitly accounting for training set overlap across repeated folds, this framework resolves the pervasive issue of pseudoreplication in standard cross-validation. This methodological contribution enables mathematically sound hypothesis testing between architectures, preventing the Type I error inflation that frequently confounds model selection in cheminformatics.
	
	\item \textbf{Architectural Benchmarking in Data-Limited Regimes:} A comprehensive empirical comparison evaluating the predictive reliability of traditional domain-engineered classical ensembles against state-of-the-art Graph Neural Networks (specifically, the ChemProp Directed Message Passing Neural Network, D-MPNN) using a standardized reference dataset.
	
	\item \textbf{Empirical Validation of Parsimony over Complexity:} Explicit demonstration that simple classical machine learning ensembles achieve statistical parity ($R^2 \approx \PaperBaselineRTwo$) with baseline deep learning models reported in recent literature. This establishes that competitive predictive reliability can be attained "out-of-the-box," effectively eliminating the necessity for synthetic data augmentation and computationally expensive 3D geometric preprocessing in low-data regimes.
	
	\item \textbf{Translational Drug Repurposing Pipeline:} The integration of machine learning predictions with structure-based molecular docking. This hybrid pipeline systematically screens, identifies, and prioritizes FDA-approved compounds, providing a validated computational pathway for rapid therapeutic repurposing.
\end{itemize}